% !TEX TS-program = xelatex
\documentclass[12pt,a4paper]{article}

\usepackage[a4paper,top=2.2cm,bottom=2.2cm,left=2.2cm,right=2.2cm]{geometry}
\usepackage{array}
\usepackage{booktabs}
\usepackage{longtable}
\usepackage{enumitem}
\usepackage{xcolor}
\usepackage{listings}
\usepackage{hyperref}
\usepackage{xepersian}

\settextfont{Yas}
\setlatintextfont{Times New Roman}
\defpersianfont\titr[Scale=1.1]{Yas}

\hypersetup{
    colorlinks=true,
    linkcolor=black,
    urlcolor=blue,
    pdftitle={گزارش نهایی پروژه برنامه‌نویسی وب},
}

\setlength{\parindent}{0pt}
\setlength{\parskip}{0.55em}
\renewcommand{\arraystretch}{1.35}

\lstset{
    basicstyle=\ttfamily\small,
    breaklines=true,
    frame=single,
    columns=fullflexible,
    showstringspaces=false,
    tabsize=2
}

\newcommand{\eng}[1]{\lr{#1}}

\begin{document}

% توجه تیم قبل از تحویل:
% تقسیم کار بخش اول بر اساس ساختار نهایی پروژه، روند توسعه و گفتگوهای تیمی تدوین شده است.
% در صورت تفاوت با تقسیم مسئولیت واقعی، فقط جدول بخش ۱ را مطابق تاریخچه واقعی تیم اصلاح کنید.

\begin{titlepage}
    \centering
    \vspace*{2cm}

    {\titr\Huge گزارش نهایی پروژه برنامه‌نویسی وب\par}
    \vspace{0.8cm}
    {\titr\Large پیاده‌سازی سرویس استریم موسیقی\par}

    \vspace{1.5cm}

    \begin{tabular}{>{\raggedleft\arraybackslash}p{7cm} p{4cm}}
        \textbf{نام و نام خانوادگی} & \textbf{شماره دانشجویی} \\ \hline
        محمدرضا ایزدی & 402110916 \\
        محمدپارسا جعفرنژادی & 402111367 \\
        سینا اصولی نژاد & 402170884 \\
    \end{tabular}

    \vfill

    \textbf{مخزن پروژه:}\\
    \begin{LTR}
    \url{https://github.com/sinaon13/web-programming-project-sut-14042}
    \end{LTR}

    \vspace{1cm}
\end{titlepage}

\tableofcontents
\newpage

% =========================================================
\section{نقش و کارهای کلی هر عضو در هر یک از دو فاز}

پروژه در دو فاز اصلی توسعه داده شد. در فاز اول تمرکز اصلی بر طراحی و پیاده‌سازی رابط کاربری، ساخت جریان‌های تعاملی، داده‌های ماک و آماده‌سازی معماری فرانت‌اند برای اتصال آینده به بک‌اند بود. در فاز دوم، بک‌اند جنگو طراحی و پیاده‌سازی شد و سپس رابط کاربری فاز اول به \eng{API} واقعی متصل گردید. با توجه به ماهیت گروهی پروژه، بخش قابل توجهی از طراحی، بازبینی، تست و رفع اشکال به شکل مشترک انجام شد و مسئولیت‌ها کاملاً منفک از یکدیگر نبودند.

\subsection{محمدرضا ایزدی}

\textbf{فاز اول:}
\begin{itemize}[nosep]
    \item مشارکت در توسعه و تست صفحات کاربری و جریان‌های تعاملی اصلی.
    \item کار روی بخش‌های مرتبط با پخش‌کننده موسیقی، پلی‌لیست‌ها و تجربه کاربری آن‌ها.
    \item بررسی رفتار صفحات در سناریوهای مختلف کاربر و کمک به رفع خطاهای رابط کاربری.
    \item مشارکت در تست دستی، هماهنگ‌سازی اجزای رابط کاربری و بررسی سناریوهای نقش‌های مختلف.
\end{itemize}

\textbf{فاز دوم:}
\begin{itemize}[nosep]
    \item تمرکز بر ادغام فرانت‌اند و بک‌اند و تست سرتاسری (\eng{End-to-End}) سناریوهای اصلی.
    \item بررسی محدودیت‌های اشتراک پایه، نقره‌ای و طلایی و کنترل رفتار سامانه پس از انقضای اشتراک.
    \item عیب‌یابی پخش فایل‌های آپلودشده و بررسی رفتار \eng{HTTP Range} برای فعال شدن صحیح \eng{Seek} در نوار زمان آهنگ.
    \item تست نقش‌های \eng{Listener}، \eng{Artist}، \eng{Support} و \eng{Admin} و بررسی سطح دسترسی‌ها در رابط کاربری و \eng{API}.
    \item مشارکت در بازبینی نسخه نهایی و جمع‌بندی مستندات پروژه.
\end{itemize}

\subsection{محمدپارسا جعفرنژادی}

\textbf{فاز اول:}
\begin{itemize}[nosep]
    \item مشارکت در توسعه صفحات تنظیمات، اعلان‌ها و پنل‌های مدیریتی/پشتیبانی.
    \item پیاده‌سازی و اصلاح اجزای رابط کاربری وابسته به نقش کاربر و وضعیت حساب هنرمند.
    \item مشارکت در واکنش‌گرا کردن صفحات و یکسان‌سازی الگوهای نمایش اطلاعات و فرم‌ها.
    \item تست جریان تیکت، تأیید/رد هنرمند و سایر سناریوهای مدیریتی در نسخه ماک.
\end{itemize}

\textbf{فاز دوم:}
\begin{itemize}[nosep]
    \item مشارکت در توسعه و بازبینی بخش‌های اشتراک (\eng{Subscription})، پشتیبانی (\eng{Support})، اعلان‌ها (\eng{Notification}) و گزارش‌گیری (\eng{Reporting}).
    \item بررسی منطق خرید اشتراک، تراکنش، قیمت‌گذاری پویا و سناریوهای انقضای اشتراک.
    \item مشارکت در پیاده‌سازی و تست جریان تیکت‌های پشتیبانی و فرایند تأیید حساب هنرمند.
    \item بازبینی گزارش‌های مدیریتی و هنرمند و اطمینان از انجام تجمیع داده‌ها (\eng{Aggregation}) در بک‌اند.
    \item مشارکت در تست رابط‌های برنامه‌نویسی (\eng{API}) و اصلاح ناسازگاری‌های قرارداد بین فرانت‌اند و بک‌اند.
\end{itemize}

\subsection{سینا اصولی نژاد}

\textbf{فاز اول:}
\begin{itemize}[nosep]
    \item مشارکت در معماری کلی پروژه فرانت‌اند، تعریف تایپ‌ها، کانتکست‌ها و ساختار صفحات.
    \item توسعه بخش‌های آرشیو موسیقی، صفحات هنرمند، آلبوم و مدیریت آثار.
    \item طراحی لایه داده ماک و آماده‌سازی ساختار کد به شکلی که جایگزینی داده‌های \eng{LocalStorage} با \eng{API} در فاز دوم کم‌هزینه باشد.
    \item مشارکت در توسعه پخش‌کننده، صف پخش، تکرار و پخش تصادفی و سایر تعاملات اصلی.
\end{itemize}

\textbf{فاز دوم:}
\begin{itemize}[nosep]
    \item مشارکت در ایجاد ساختار اصلی پروژه \eng{Django} و تفکیک آن به اپ‌های مستقل.
    \item توسعه و بازبینی مدل‌ها و رابط‌های برنامه‌نویسی بخش‌های \eng{Accounts}، \eng{Music} و \eng{Playlists} و اتصال آن‌ها به فرانت‌اند.
    \item پیاده‌سازی/بازبینی احراز هویت \eng{JWT}، پرمیشن‌ها، سریالایزرها، ویوست‌ها و اندپوینت‌های اصلی.
    \item مشارکت در مدیریت فایل‌های صوتی و تصاویر، مستندسازی \eng{API} و تنظیمات پروژه.
    \item مشارکت در تست‌های بک‌اند، رفع اشکال‌های نهایی و آماده‌سازی اجرای پروژه با Docker.
\end{itemize}

% =========================================================
\section{قوانین توسعه و فرآیند آن}

برای کاهش ناسازگاری بین بخش‌های مختلف و افزایش قابلیت نگهداری، مجموعه‌ای از قواعد توسعه در طول پروژه رعایت شد.

\subsection{قواعد نام‌گذاری}

\begin{itemize}
    \item در کد TypeScript/React، نام متغیرها و توابع با \eng{camelCase} نوشته شدند؛ مانند
    \eng{currentUser}، \eng{playTrack} و \eng{subscriptionDaysLeft}.
    \item نام Componentها، Typeها و Contextها با \eng{PascalCase} تعریف شدند؛ مانند
    \eng{PlayerContext}، \eng{AuthContext} و \eng{DownloadButton}.
    \item در بک‌اند Python/Django از قرارداد \eng{snake\_case} برای متغیرها، توابع و فیلدهای مدل استفاده شد؛ مانند
    \eng{expires\_at}، \eng{audio\_file} و \eng{notification\_type}.
    \item نام کلاس‌های Django با \eng{PascalCase} تعریف شدند؛ مانند
    \eng{CustomUser}، \eng{SubscriptionPlan} و \eng{TrackStreamView}.
    \item نام نشانی‌ها و مسیرها کوتاه، قابل پیش‌بینی و همسو با اصول \eng{REST} نگه داشته شد.
\end{itemize}

\subsection{قواعد \eng{Git} و شاخه‌ها}

شاخه اصلی نسخه پایدار پروژه \eng{main} در نظر گرفته شد. تغییرات توسعه‌ای به صورت موضوعی و با کامیت‌های کوچک‌تر و قابل ردیابی اعمال شدند. برای نام‌گذاری شاخه‌های موقت در توسعه گروهی از الگوی زیر استفاده شد:

\begin{LTR}
\begin{verbatim}
feature/<feature-name>
fix/<bug-name>
refactor/<module-name>
\end{verbatim}
\end{LTR}

پیام کامیت‌ها نیز تا حد امکان هدف تغییر را مشخص می‌کردند؛ برای مثال کامیت‌های مرتبط با اصلاح پخش‌کننده، امنیت بک‌اند، پلی‌لیست و اشتراک به صورت مستقل ثبت شدند. در پایان هر مرحله، تغییرات تست و سپس در شاخه اصلی تثبیت شدند.

\subsection{فرآیند توسعه}

فرآیند کلی توسعه در هر قابلیت به شکل زیر بود:

\begin{enumerate}
    \item استخراج نیازمندی از مستند پروژه.
    \item طراحی داده و تعیین مسئولیت فرانت‌اند و بک‌اند.
    \item پیاده‌سازی اولیه در یک بخش محدود از سیستم.
    \item تست دستی سناریوی عادی و سناریوهای خطا.
    \item در فاز دوم، اتصال قابلیت به \eng{API} و حذف وابستگی قابلیت اصلی به داده ماک.
    \item بررسی سطح دسترسی و اعتبارسنجی در سمت بک‌اند؛ به‌خصوص در قابلیت‌های مالی، اشتراک، آثار هنرمندان و پنل مدیریتی.
    \item اجرای تست‌های خودکار و اصلاح خطاهای بازگشتی (\eng{Regression}).
    \item بازبینی کد از نظر تکرار، نام‌گذاری و قرارداد درخواست و پاسخ.
\end{enumerate}

در طراحی \eng{API} تلاش شد امنیت و منطق تجاری تنها به رابط کاربری وابسته نباشد و بررسی‌های حساس در بک‌اند انجام شوند. همچنین از \eng{Swagger/OpenAPI} برای مشاهده و تست قرارداد اندپوینت‌ها استفاده شد.

% =========================================================
\section{ساختار هر یک از دو پروژه و توجیه آن}

با وجود نگهداری کل پروژه در یک مخزن (\eng{Repository})، فرانت‌اند و بک‌اند از نظر منطقی دو پروژه مستقل هستند و هرکدام ساختار و چرخه اجرای مخصوص خود را دارند.

\subsection{ساختار فرانت‌اند}

فرانت‌اند با \eng{Next.js}، \eng{React} و \eng{TypeScript} توسعه داده شده است. ساختار اصلی آن به صورت زیر است:

\begin{LTR}
\begin{verbatim}
src/
  app/
    admin/
    albums/
    artist/
    artist-portal/
    browse/
    login/
    notifications/
    playlists/
    profile/
    register/
    settings/
    singles/
    support-tickets/
  components/
    layout/
    player/
    pwa/
    ui/
  context/
    AuthContext.tsx
    LanguageContext.tsx
    PlayerContext.tsx
  lib/
    api.ts
    adapters.ts
    types.ts
    mockData.ts
  __tests__/
public/
__tests__/
\end{verbatim}
\end{LTR}

\textbf{توجیه ساختار فرانت‌اند:}

\begin{itemize}
    \item پوشه \eng{app} مسئول مسیرها و صفحات است و با معماری \eng{App Router} در \eng{Next.js} هماهنگ است.
    \item کامپوننت‌های مشترک از صفحات جدا شده‌اند تا کد تکراری کمتر شود.
    \item کانتکست‌ها مسئول وضعیت‌های سراسری مانند احراز هویت، زبان و پخش‌کننده هستند؛ بنابراین نیاز به انتقال زنجیره‌ای پراپ‌ها (\eng{Props}) بین صفحات کاهش یافته است.
    \item فایل \eng{api.ts} نقطه مرکزی ارتباط با بک‌اند است و مدیریت توکن، نوسازی توکن (\eng{Refresh}) و اندپوینت‌ها را یکجا انجام می‌دهد.
    \item فایل \eng{adapters.ts} مرز بین قرارداد داده \eng{Django} و تایپ‌های فرانت‌اند است و تبدیل \eng{snake\_case} به ساختار مورد نیاز رابط کاربری را انجام می‌دهد.
    \item تایپ‌ها در یک فایل مستقل قرار گرفته‌اند تا قرارداد داخلی فرانت‌اند مشخص و قابل بررسی باشد.
    \item داده‌های ماک فاز اول در لایه مستقل نگهداری شدند تا در فاز دوم امکان جایگزینی آن‌ها با \eng{API} واقعی بدون بازنویسی کل رابط کاربری فراهم شود.
\end{itemize}

\subsection{ساختار بک‌اند}

بک‌اند با \eng{Django} و \eng{Django REST Framework} توسعه داده شده و به چند اپ مستقل تقسیم شده است:

\begin{LTR}
\begin{verbatim}
backend/
  accounts/
  music/
  playlists/
  subscriptions/
  notifications/
  support/
  reports/
  config/
  manage.py
  pyproject.toml
  uv.lock
\end{verbatim}
\end{LTR}

هر اپ در صورت نیاز شامل اجزای معمول زیر است:

\begin{LTR}
\begin{verbatim}
models.py
serializers.py
views.py
urls.py
permissions.py
tests.py
migrations/
\end{verbatim}
\end{LTR}

\textbf{توجیه ساختار بک‌اند:}

\begin{itemize}
    \item هر اپ یک حوزه کسب‌وکار مستقل را مدیریت می‌کند؛ برای مثال اپ \eng{Music} مسئول آثار و استریم‌ها و اپ \eng{Subscriptions} مسئول پلن‌ها و تراکنش‌ها است.
    \item این تفکیک باعث کاهش وابستگی متقابل (\eng{Coupling}) و ساده‌تر شدن توسعه و تست مستقل هر بخش می‌شود.
    \item سریالایزرها مسئول اعتبارسنجی و تبدیل داده هستند و ویوها مسئول رفتار \eng{HTTP} و کنترل دسترسی‌اند؛ در نتیجه مسئولیت‌ها با هم مخلوط نمی‌شوند.
    \item پرمیشن‌های مشترک برای نقش‌ها به صورت کلاس‌های مستقل تعریف شدند.
    \item از ویوهای عمومی و ویوست‌های \eng{DRF} برای کاهش کد تکراری عملیات \eng{CRUD} استفاده شده است.
    \item بخش \eng{Reports} به جای ذخیره داده تکراری، آمار را از مدل‌های اصلی به‌صورت تجمیعی (\eng{Aggregation}) محاسبه می‌کند.
    \item فایل‌های مایگریشن تاریخچه ساختار پایگاه داده را نگه می‌دارند و امکان اجرای قابل تکرار پروژه را فراهم می‌کنند.
\end{itemize}

\subsection{فناوری‌های اصلی}

در فرانت‌اند از \eng{Next.js}، \eng{React}، \eng{TypeScript}، \eng{Tailwind CSS} و \eng{Recharts} استفاده شده است. در بک‌اند از \eng{Django}، \eng{Django REST Framework}، \eng{SimpleJWT}، \eng{django-filter}، \eng{django-cors-headers}، \eng{drf-spectacular} و \eng{Pillow} استفاده شد. پایگاه داده نسخه توسعه \eng{SQLite} است. برای اجرای قابل حمل نیز \eng{Dockerfile} و \eng{docker-compose} در مخزن پروژه قرار داده شدند.

% =========================================================
\section{توجیه نگهداری‌پذیری کد}

نگهداری‌پذیری یکی از اهداف معماری پروژه بود. تصمیم‌های زیر در این زمینه گرفته شدند:

\subsection{تفکیک مسئولیت‌ها}

منطق نمایش در کامپوننت‌ها، وضعیت‌های سراسری در کانتکست‌ها، ارتباط شبکه در کلاینت \eng{API} و تبدیل قرارداد داده در آداپترها از یکدیگر جدا شده‌اند. در بک‌اند نیز مدل، سریالایزر، ویو و پرمیشن مسئولیت‌های مشخصی دارند. این جداسازی باعث می‌شود تغییر در یک بخش، کمترین اثر جانبی را روی بخش‌های دیگر داشته باشد.

\subsection{استفاده از قراردادهای مشخص داده}

وجود تایپ‌های \eng{TypeScript} در فرانت‌اند و سریالایزرهای \eng{DRF} در بک‌اند، قرارداد داده را شفاف می‌کند. استفاده از آداپترها نیز اجازه می‌دهد تغییر نام یا ساختار فیلدهای بک‌اند لزوماً به بازنویسی همه کامپوننت‌ها منجر نشود.

\subsection{استفاده از امکانات استاندارد فریم‌ورک}

به جای پیاده‌سازی دستی عملیات تکراری، از \eng{ViewSet}، \eng{GenericAPIView}، سریالایزر، پرمیشن و فیلترهای \eng{DRF} استفاده شده است. این موضوع حجم کد سفارشی را کاهش داده و رفتار سیستم را قابل پیش‌بینی‌تر می‌کند.

\subsection{مدیریت متمرکز احراز هویت و دسترسی}

توکن JWT و Refresh Token در یک API Client مرکزی مدیریت می‌شوند. در سمت بک‌اند نیز نقش‌ها و دسترسی‌های حساس با Permissionهای مشخص کنترل می‌شوند. در نتیجه قواعد امنیتی در هر صفحه به شکل مستقل و متفاوت تکرار نمی‌شوند.

\subsection{تست‌پذیری}

برای اپ‌های مختلف بک‌اند تست‌های مستقل نوشته شده است و سناریوهایی مانند ثبت‌نام، نقش‌ها، آثار، پلی‌لیست، اشتراک، اعلان، پشتیبانی و گزارش‌گیری پوشش داده شده‌اند. در فرانت‌اند نیز تست‌هایی برای کانتکست، پخش‌کننده و کامپوننت‌ها وجود دارد. وجود تست‌ها ریسک خطاهای بازگشتی هنگام بازآرایی (\eng{\eng{Refactor}}) را کاهش می‌دهد.

\subsection{مستندسازی و مشاهده API}

استفاده از \eng{drf-spectacular} باعث تولید شِما و رابط \eng{Swagger UI} شده است. این ابزار برای بررسی سریع اندپوینت‌ها و درخواست/پاسخ‌ها و اتصال فرانت‌اند و بک‌اند مفید است.

\subsection{قابلیت توسعه آینده}

تفکیک حوزه‌ها در اپ‌های \eng{Django} و مستقل بودن لایه \eng{API} در فرانت‌اند اجازه می‌دهد قابلیت‌هایی مانند سیستم پیشنهاددهنده، سیستم پرداخت کامل‌تر، فضای ذخیره‌سازی خارجی یا پایگاه داده محیط عملیاتی (\eng{Production}) بدون تغییر بنیادی کل پروژه اضافه شوند.

% =========================================================
\section{مدل‌های موجود در بک‌اند، روابط آن‌ها و توجیه وجودشان}

\subsection{مدل‌های Accounts}

\subsubsection{\eng{CustomUser}}

مدل اصلی کاربر است و از \eng{AbstractUser} توسعه داده شده است. علاوه بر اطلاعات پایه، فیلدهایی برای نام نمایشی، نقش، تصویر، تاریخ تولد، جنسیت، زندگی‌نامه، اطلاعات درخواست هنرمند و وضعیت درآمد هنرمند دارد.

نقش‌های سیستم در این مدل عبارت‌اند از:
\eng{LISTENER}، \eng{ARTIST}، \eng{SUPPORT} و \eng{ADMIN}.

همچنین رابطه \eng{Following} به صورت \eng{ManyToMany} خودارجاعی و غیرمتقارن تعریف شده است؛ بنابراین اگر کاربر «الف» کاربر «ب» را دنبال کند، الزاماً کاربر «ب» دنبال‌کننده کاربر «الف» نیست.

\textbf{توجیه:}
به دلیل وجود چهار نقش و اطلاعات مشترک فراوان، استفاده از یک مدل کاربر (\eng{User Model}) توسعه‌یافته به جای چند مدل جداگانه، احراز هویت و پرمیشن‌ها را ساده‌تر و یکپارچه‌تر می‌کند.

\subsubsection{\eng{UserPreferences}}

رابطه آن با کاربر از نوع \eng{OneToOne} است و تنظیماتی مانند زبان، حجم صدای پیش‌فرض و فعال بودن اعلان‌ها را ذخیره می‌کند.

\textbf{توجیه:}
تنظیمات باید بین دستگاه‌های مختلف همگام شوند؛ بنابراین نگهداری آن‌ها در بک‌اند به جای وابستگی کامل به \eng{LocalStorage} ضروری است.

\subsection{مدل‌های Music}

\subsubsection{\eng{Album}}

هر آلبوم با یک \eng{ForeignKey} به هنرمند متصل است و اطلاعاتی مانند عنوان، کاور، تاریخ انتشار و ژانر را ذخیره می‌کند.

\textbf{توجیه:}
آلبوم موجودیتی مستقل است که چند ترک را گروه‌بندی می‌کند و اطلاعات مشترک انتشار را نگه می‌دارد.

\subsubsection{\eng{Track}}

هر ترک با یک هنرمند رابطه \eng{ForeignKey} دارد و می‌تواند به صورت اختیاری به یک آلبوم متصل باشد. اطلاعات فایل اصلی صوتی، نسخه کیفیت پایین، کاور، متن، ژانر، نوع انتشار، همکاران، فرمت و دسترسی زودهنگام (\eng{Early Access}) در این مدل قرار گرفته‌اند.

\textbf{توجیه:}
ترک هسته اصلی سرویس استریم است و باید هم فایل رسانه‌ای و هم فراداده مورد نیاز پخش‌کننده و جست‌وجو، گزارش‌گیری و محدودیت دسترسی را در اختیار سیستم قرار دهد.

\subsubsection{\eng{StreamLog}}

هر \eng{StreamLog} با یک کاربر و یک ترک دو رابطه \eng{ForeignKey} دارد و زمان استریم را ثبت می‌کند.

\textbf{توجیه:}
این مدل برای محدودیت روزانه کاربران پایه (\eng{Basic})، محاسبه تعداد استریم و شنونده یکتا (\eng{Unique Listener})، گزارش‌گیری و محاسبات مالی ضروری است. ذخیره لاگ مستقل باعث می‌شود گزارش‌ها قابل محاسبه و قابل بازسازی باشند.

\subsection{مدل‌های Playlists}

\subsubsection{\eng{Playlist}}

هر پلی‌لیست به مالک خود از طریق \eng{ForeignKey} متصل است و عنوان، توضیح، کاور، عمومی/خصوصی بودن و زمان‌های ایجاد و ویرایش را ذخیره می‌کند.

\subsubsection{\eng{PlaylistTrack}}

مدل واسط بین پلی‌لیست و ترک است و علاوه بر رابطه دو طرف، موقعیت و زمان افزودن آهنگ را نیز نگه می‌دارد. زوج پلی‌لیست و ترک یکتا در نظر گرفته شده است.

\textbf{توجیه:}
استفاده از مدل واسط نسبت به یک رابطه ساده \eng{ManyToMany} ساده امکان نگهداری ترتیب آهنگ‌ها و فراداده مربوط به عضویت آهنگ در پلی‌لیست را فراهم می‌کند.

\subsection{مدل‌های Subscriptions}

\subsubsection{\eng{SubscriptionPlan}}

اطلاعات هر سطح اشتراک شامل سطح اشتراک، قیمت، مدت، محدودیت روزانه استریم، حداکثر تعداد پلی‌لیست، توضیح و فعال بودن پلن را ذخیره می‌کند.

\textbf{توجیه:}
قیمت‌گذاری و محدودیت‌ها نباید به‌صورت ثابت در کد (\eng{Hard-coded}) شوند. با وجود این مدل، مدیر می‌تواند قیمت را در پایگاه داده تغییر دهد و Backend همان مقدار جدید را استفاده کند.

\subsubsection{\eng{UserSubscription}}

رابطه آن با کاربر از نوع \eng{OneToOne} و با پلن از نوع \eng{ForeignKey} است. زمان شروع، زمان انقضا، فعال بودن و تمدید خودکار (\eng{Auto Renew}) را نگه می‌دارد.

\textbf{توجیه:}
برای تشخیص سطح اشتراک فعلی کاربر و کنترل انقضای اشتراک نیاز به یک وضعیت جاری و قابل پرس‌وجو وجود دارد.

\subsubsection{\eng{Transaction}}

هر تراکنش به یک کاربر و یک \eng{SubscriptionPlan} متصل است و اطلاعات ماه‌های خریداری‌شده، مبلغ، \eng{Authority}، شناسه مرجع (\eng{Ref ID})، وضعیت تراکنش و زمان تأیید را نگه می‌دارد.

\textbf{توجیه:}
طبق نیازمندی پروژه، تراکنش‌ها باید در وضعیت‌های در حال انجام، موفق یا ناموفق قابل پیگیری باشند و سابقه پرداخت از اشتراک جاری جدا بماند.

\subsection{مدل Notifications}

\subsubsection{\eng{Notification}}

هر اعلان به یک دریافت‌کننده (\eng{Recipient}) متصل است و عنوان، پیام، نوع، وضعیت خوانده‌شدن، لینک مقصد و زمان ایجاد را ذخیره می‌کند.

\textbf{توجیه:}
یک مدل عمومی اعلان اجازه می‌دهد رویدادهای مختلف مانند اشتراک، انتشار موسیقی، دنبال‌کردن، پشتیبانی و پرداخت با ساختار یکسان در رابط کاربری نمایش داده شوند.

\subsection{مدل‌های Support}

\subsubsection{\eng{Ticket}}

هر تیکت یک کاربر ایجادکننده دارد و می‌تواند به یک کاربر پشتیبان یا مدیر از طریق \eng{assigned\_to} متصل شود. موضوع، توضیحات، وضعیت و اولویت و زمان‌ها در آن نگهداری می‌شوند.

\subsubsection{\eng{TicketMessage}}

هر پیام به یک تیکت و یک فرستنده متصل است و متن و زمان پیام را ذخیره می‌کند.

\textbf{توجیه:}
تفکیک تیکت از پیام باعث می‌شود یک تیکت بتواند رشته گفت‌وگوی چندپیامی (\eng{Thread}) داشته باشد و هم کاربر و هم پشتیبان در همان جریان گفتگو مشارکت کنند.

\subsection{بخش \eng{Reports}}

بخش \eng{Reports} مدل دائمی جداگانه‌ای ایجاد نکرده است و داده‌های آماری را از مدل‌های موجود مانند \eng{User}، \eng{Subscription}، \eng{Transaction}، \eng{Track} و \eng{StreamLog} با کوئری‌های تجمیعی (\eng{Aggregate}) محاسبه می‌کند.

\textbf{توجیه:}
بخش زیادی از داده‌های گزارش، داده مشتق‌شده هستند. ذخیره دوباره آن‌ها می‌تواند باعث ناسازگاری شود؛ بنابراین محاسبه تجمیع داده‌ها در بک‌اند هم با نیازمندی پروژه هماهنگ است و هم از تکرار داده جلوگیری می‌کند.

% =========================================================

\section{نقش هوش مصنوعی در توسعه}

در این پروژه، مطابق تأکیدی که در جلسات درس نیز مطرح شد، هوش مصنوعی به عنوان
\textbf{دستیار تسهیل‌کننده و شتاب‌دهنده توسعه} مورد استفاده قرار گرفت، نه به عنوان جایگزین توسعه‌دهنده.
ابزارهای هوش مصنوعی در بخش‌های مختلف می‌توانند سرعت تولید کد و تحلیل مسئله را افزایش دهند، اما مسئولیت نهایی
صحت، امنیت و سازگاری کد با منطق پروژه بر عهده اعضای تیم باقی می‌ماند.

\subsection{کاربردهای اصلی هوش مصنوعی در پروژه}

\begin{itemize}
    \item \textbf{تولید ساختار اولیه و کدهای تکراری:}
    از هوش مصنوعی برای تولید نسخه اولیه بخش‌هایی مانند مدل‌های پایگاه داده، سریالایزرها،
    ویوها، تست‌ها، ساختار کامپوننت‌ها و کدهای تکراری استفاده شد. این کار باعث شد زمان بیشتری
    برای بررسی منطق اصلی پروژه و تست سناریوها صرف شود.

    \item \textbf{کاهش خطاهای کوچک برنامه‌نویسی:}
    ابزارهای هوش مصنوعی در بازبینی کد و شناسایی خطاهای ساده مانند اشتباه‌های تایپی،
    ناسازگاری نام فیلدها، خطاهای ناشی از کپی‌کردن کد و مشکلات جزئی در قراردادهای داده مفید بودند.

    \item \textbf{کمک در پیکربندی و اشکال‌زدایی:}
    در زمان توسعه برای تحلیل خطاهای \eng{Django} و \eng{Next.js}، تنظیم
    \eng{CORS}، \eng{JWT}، \eng{Swagger/OpenAPI}، وابستگی‌های پروژه، متغیرهای محیطی
    و مشکلات ارتباط فرانت‌اند و بک‌اند از هوش مصنوعی کمک گرفته شد.

    \item \textbf{کمک در نوشتن کوئری‌ها و منطق داده:}
    برای طراحی بعضی کوئری‌های پیچیده \eng{ORM}، محاسبه داده‌های تجمیعی و بررسی روابط
    بین مدل‌ها از مدل‌های زبانی استفاده شد. خروجی پیشنهادی پس از آن با ساختار واقعی مدل‌ها و
    نیازمندی پروژه تطبیق داده شد.

    \item \textbf{تحلیل نیازمندی و طراحی تست:}
    از هوش مصنوعی برای تبدیل مستند پروژه به سناریوهای قابل تست، به‌خصوص در زمینه نقش‌ها،
    سطح اشتراک، انقضای اشتراک، دسترسی زودهنگام، دانلود، تعداد پلی‌لیست‌ها و محدودیت استریم استفاده شد.

    \item \textbf{کمک در مستندسازی:}
    در جمع‌بندی معماری، توضیح مدل‌ها، روابط، تصمیم‌های طراحی و آماده‌سازی گزارش نهایی نیز
    از هوش مصنوعی به عنوان ابزار کمکی استفاده شد.
\end{itemize}

\subsection{محدودیت‌ها و مسئولیت توسعه‌دهنده}

یکی از نکات مهم مطرح‌شده در جلسات درس این بود که
\textbf{مسئولیت پروژه را نمی‌توان به ابزار هوش مصنوعی منتقل کرد}.
مدل زبانی ممکن است کدی تولید کند که از نظر نحوی صحیح باشد، اما با منطق تجاری پروژه سازگار نباشد.
برای مثال، تعیین الزامی یا اختیاری بودن یک فیلد، طول مجاز داده، سطح دسترسی یک نقش، نحوه تمدید اشتراک
یا شرایط اعتبار یک تراکنش، نیازمند شناخت دقیق زمینه و منطق پروژه است.

به همین دلیل در این پروژه خروجی هوش مصنوعی بدون بررسی مستقیم وارد نسخه نهایی نشد. کدهای پیشنهادی
با اجرای واقعی برنامه، تست \eng{API}، مشاهده رفتار مرورگر، تست‌های خودکار و بازبینی انسانی بررسی شدند.
این موضوع به‌خصوص در بخش‌های مالی، احراز هویت، سطح دسترسی، اشتراک و فایل‌های صوتی اهمیت بیشتری داشت.

همچنین بین استفاده آموزشی و توسعه صنعتی تفاوت قائل شدیم. در محیط آموزشی ممکن است برای سرعت بیشتر
یک قطعه‌کد اولیه از هوش مصنوعی دریافت شود، اما برای رسیدن به نسخه قابل اتکا باید جزئیات ورودی، خروجی،
حالت‌های خطا، محدودیت فیلدها و رفتار سناریوهای مرزی به صورت صریح بررسی و سفارشی‌سازی شوند.

\section{یک نمونه از کد هوش مصنوعی حاضر در هر یک از دو فاز}

\subsection{نمونه فاز اول: لایه ساده داده ماک در فرانت‌اند}

در فاز اول بک‌اند واقعی وجود نداشت و بخشی از داده‌ها در \eng{LocalStorage} نگهداری می‌شدند.
یکی از قطعه‌کدهایی که با کمک هوش مصنوعی طراحی شد، تابع‌های عمومی خواندن و نوشتن داده‌های ماک بود:

\begin{LTR}
\begin{lstlisting}
export const getDB = <T>(key: string, fallback: T): T => {
  if (typeof window === 'undefined') return fallback;

  const item = localStorage.getItem(key);
  return item ? JSON.parse(item) : fallback;
};

export const setDB = <T>(key: string, data: T): void => {
  if (typeof window !== 'undefined') {
    localStorage.setItem(key, JSON.stringify(data));
  }
};
\end{lstlisting}
\end{LTR}

این کد نمونه‌ای از استفاده مناسب هوش مصنوعی برای تولید
\textbf{کد تکراری و ساختار اولیه} است. پس از تولید اولیه، نحوه استفاده از آن با ساختار واقعی پروژه هماهنگ شد
تا کامپوننت‌ها مستقیماً به جزئیات ذخیره‌سازی وابسته نباشند.

\subsection{نمونه فاز دوم: پشتیبانی \eng{HTTP Range} برای فایل صوتی}

در فاز دوم هنگام تست ترک‌های آپلودشده مشخص شد که مرورگر برای جابه‌جایی در زمان آهنگ
از درخواست محدوده‌ای (\eng{Range Request}) استفاده می‌کند. با کمک هوش مصنوعی ابتدا رفتار درخواست و پاسخ
در بخش شبکه مرورگر بررسی شد و سپس راه‌حل زیر برای پاسخ \eng{206 Partial Content} طراحی شد:

\begin{LTR}
\begin{lstlisting}[language=Python]
def serve_media_with_range(request, path):
    file_path = os.path.join(settings.MEDIA_ROOT, path)

    if not os.path.exists(file_path):
        raise Http404("Media file not found")

    file_size = os.path.getsize(file_path)
    range_header = request.META.get("HTTP_RANGE", "").strip()

    if range_header:
        range_match = re.search(r"bytes=(\d+)-(\d*)", range_header)

        if range_match:
            first_byte = int(range_match.group(1))
            last_byte = (
                int(range_match.group(2))
                if range_match.group(2)
                else file_size - 1
            )

            length = last_byte - first_byte + 1

            response = StreamingHttpResponse(
                file_iterator(),
                status=206,
                content_type=content_type,
            )

            response["Content-Range"] = (
                f"bytes {first_byte}-{last_byte}/{file_size}"
            )
            response["Accept-Ranges"] = "bytes"
            response["Content-Length"] = str(length)

            return response
\end{lstlisting}
\end{LTR}

در این نمونه، نقش هوش مصنوعی تنها «تولید کد» نبود؛ ابتدا از آن برای تحلیل علت خطا استفاده شد.
پس از مشاهده تفاوت پاسخ \eng{200} و \eng{206} در مرورگر، راه‌حل با رفتار واقعی فایل صوتی تطبیق داده شد
و سپس نتیجه به صورت دستی تست شد.

\section{ضعف‌ها و قوت‌های هوش مصنوعی در توسعه بک‌اند و فرانت‌اند}

در جلسات درس، هوش مصنوعی به صورت یک ابزار عمومی برای هر دو بخش فرانت‌اند و بک‌اند معرفی و استفاده شد
و مقایسه صریحی که یکی از این دو حوزه را ذاتاً مناسب‌تر یا نامناسب‌تر برای هوش مصنوعی بداند ارائه نشد.
بنابراین موارد زیر ترکیبی از چارچوب مطرح‌شده در کلاس و تجربه عملی تیم در همین پروژه است.

\subsection{فرانت‌اند}

\subsubsection{قوت‌ها}

\begin{itemize}
    \item تولید سریع ساختار اولیه صفحات، کامپوننت‌ها، فرم‌ها و استایل‌های تکراری.
    \item کمک در طراحی کانتکست‌ها، مدیریت وضعیت و اتصال کامپوننت‌ها به \eng{API}.
    \item سرعت بالا در پیشنهاد کدهای \eng{CSS/Tailwind} و چیدمان واکنش‌گرا.
    \item کمک در پیدا کردن خطاهای ساده مربوط به نام فیلدها، تایپ‌ها و قرارداد داده.
    \item مفید بودن در اشکال‌زدایی خطاهای مرورگر، پخش‌کننده، شبکه و رفتار فایل صوتی.
    \item کمک در ساخت تست‌های رابط کاربری و سناریوهای تعاملی.
\end{itemize}

\subsubsection{ضعف‌ها}

\begin{itemize}
    \item هوش مصنوعی ممکن است نام فیلدهای فرانت‌اند و بک‌اند را ناسازگار تولید کند؛
    برای مثال یک سمت \eng{name} و سمت دیگر \eng{title}.
    \item در کدهای طولانی \eng{React} ممکن است وضعیت قدیمی، وابستگی ناقص در هوک یا بازآرایی ناقص ایجاد شود.
    \item پیشنهاد هوش مصنوعی گاهی فقط محدودیت را در رابط کاربری اعمال می‌کند؛ در حالی که دسترسی حساس باید در بک‌اند نیز اعمال شود.
    \item رفتار واقعی \eng{Media API} و مرورگر همیشه از روی کد قابل پیش‌بینی نیست و نیازمند تست با \eng{DevTools} است.
    \item کد تولیدشده ممکن است با ظاهر پروژه سازگار باشد اما زمینه دقیق نیازمندی را در نظر نگیرد.
\end{itemize}

\subsection{بک‌اند}

\subsubsection{قوت‌ها}

\begin{itemize}
    \item تولید سریع مدل‌ها، سریالایزرها، ویوها و کدهای \eng{CRUD} متداول.
    \item کمک در طراحی اعتبارسنجی، پرمیشن و روابط بین مدل‌ها.
    \item کمک در پیکربندی \eng{Swagger/OpenAPI}، \eng{CORS} و احراز هویت \eng{JWT}.
    \item کمک در تولید کوئری‌های پیچیده \eng{ORM} و داده‌های تجمیعی برای گزارش‌ها.
    \item تحلیل سریع ردگیری خطاها، خطاهای وابستگی و مشکلات اجرای پروژه.
    \item کمک در طراحی تست برای سناریوهای سطح دسترسی مبتنی بر نقش و اشتراک.
\end{itemize}

\subsubsection{ضعف‌ها}

\begin{itemize}
    \item هوش مصنوعی زمینه کامل منطق تجاری پروژه را به صورت خودکار نمی‌داند؛
    بنابراین تعیین الزامی یا اختیاری بودن فیلدها، محدودیت‌ها و رفتار دقیق سناریوها باید توسط توسعه‌دهنده انجام شود.
    \item در بخش‌های مالی و امنیتی، یک شرط اشتباه می‌تواند اثر جدی داشته باشد؛
    بنابراین کد پیشنهادی هوش مصنوعی باید حتماً بازبینی شود.
    \item ممکن است منطق مشابه در چند ویو به شکل متفاوت تولید شود و در نتیجه رفتار بخش‌های مختلف یکسان نباشد.
    \item در درگاه پرداخت، اعتبارسنجی تراکنش، تکرارپذیری امن و مدیریت خطاهای شبکه نمی‌توان تنها به کد پیشنهادی هوش مصنوعی اکتفا کرد.
    \item ممکن است یک کتابخانه در کد استفاده شود اما به فایل وابستگی پروژه اضافه نشود؛
    بنابراین اجرای پروژه در محیط تازه باید جداگانه بررسی شود.
    \item در کوئری‌های پیچیده، اگر مدل‌ها و روابط به طور کامل در پرامپت توضیح داده نشوند،
    هوش مصنوعی ممکن است کوئری ظاهراً معتبر اما از نظر منطقی اشتباه تولید کند.
\end{itemize}

\subsection{جمع‌بندی تجربه استفاده از هوش مصنوعی}

نتیجه تجربه تیم این بود که هوش مصنوعی بیشترین ارزش را زمانی ایجاد می‌کند که به عنوان
\textbf{ابزار افزایش سرعت، تولید پیش‌نویس و کمک به تحلیل} استفاده شود و نه به عنوان منبع نهایی تصمیم.
در فرانت‌اند، این مزیت بیشتر در تولید سریع رابط کاربری، کامپوننت و اشکال‌زدایی تعاملات دیده شد؛
در بک‌اند نیز در تولید ساختارهای استاندارد \eng{DRF}، اعتبارسنجی، کوئری و تحلیل خطا بسیار مفید بود.

در هر دو بخش، مرحله نهایی باید شامل خواندن و فهم کد، تطبیق با مستند پروژه، تست حالت‌های مختلف
و مسئولیت‌پذیری توسعه‌دهنده باشد. در نتیجه، استفاده از هوش مصنوعی سرعت تیم را افزایش داد،
اما تصمیم نهایی درباره معماری، منطق تجاری، امنیت و صحت رفتار سامانه توسط اعضای تیم گرفته شد.


\end{document}
